
\begin{lemma}{7.1.55}\label{lem:7-1-55}
若
$\Gamma\ne
\Gamma_0^{\mathrm f},
\Gamma_0^{\mathrm m},
\Gamma_0^{\mathrm m'}$，
则
$\mathcal N(\Gamma)$
微分同胚于
$D^{|\Gamma|}$。
\end{lemma}

事实上，
引理~\ref{lem:7-1-55}
的结论对上述三个
$\Gamma_0$ 也成立；
这些情形将在
定理~\ref{thm:7-1-51}
的证明末尾处理。
下一个引理亦有类似说明。

\begin{proof}
$\mathcal N(\Gamma)$
是与
$\overline{\mathbf S}$
各分支对应的模空间的乘积；
这里已将
$\overline{\mathbf S}$
同 $\Gamma$ 相识别。
若 $\overline{\mathbf S}_i$
是 $\mathrm m'$ 型或 $\mathrm m$ 型分支，
则相应模空间为
$\mathcal M^{\mathrm{main}}_{k_i}
\cong D^{k_i-3}$；
$k_i$ 是该分支上的奇点或标记点数。
若它是 $\mathrm f$ 型，
相应模空间为
$\mathcal N_{k_i'+1}$，
其中 $k_i'<k$，
并且 $k_i'$
等于
$\overline{\mathbf S}_i$
上的标记点数与集合
\[
\bigcup_{j\ne i}
\bigl(
\overline{\mathbf S}_i\cap
\overline{\mathbf S}_j
\bigr)
\]
的元素数之和。
所以，由归纳假设，
$\mathcal N(\Gamma)$
同胚于一个圆盘\translatornote{原书引理 7.1.55
的陈述为“微分同胚”，但证明在此只写
“同胚于一个圆盘”，并略去维数计算；
尚未交代乘积光滑结构与该层闭包的相容性，
见疑误 A-0427。}。
维数计算直接，故从略。
\end{proof}

\begin{lemma}{7.1.56}\label{lem:7-1-56}
除去那些
$\mathbf S=(\Sigma,(z_0,\ldots,z_k),\rho)$
的组合型
$\overline{\mathbf S}$
仅有一个内部顶点、且 $\rho$ 为常数的点，
$\mathcal N_{k+1}$
具有带边界与角的光滑流形结构。
\end{lemma}

\begin{proof}
取一个不属于上述情形的点
$\mathbf S
=(\Sigma,(z_0,\ldots,z_k),\rho)$。
我们将在其邻域内构造坐标图。
设 $\mathbf S$ 对应于 $\Gamma$，
并取前述
$\Sigma_{\mathrm f,i}$、
$\Sigma_{\mathrm m,i}$、
$\Sigma_{\mathrm m'}$。
如下重新分解
$\Sigma_{\mathrm f,i}$。
若其中两个分支
$\Sigma_a,\Sigma_b$
相交且
$\rho_a=\rho_b$，
则令
$\Sigma_a\sim\Sigma_b$。
在
$\Sigma_{\mathrm f,i}$
的分支集上取此关系生成的等价关系，
把同一等价类的分支合并，
得到分解
\[
\Sigma_{\mathrm f,i}
=\bigcup_{j=1}^{a_i}
\Sigma_{\mathrm f,i,j}.
\]

令
$k_{\mathrm m,i}+1$
等于
$\Sigma_{\mathrm m,i}$
上的标记点数与集合
\[
\bigcup_{*\ne(\mathrm m,i)}
\bigl(
\Sigma_{\mathrm m,i}\cap\Sigma_*
\bigr)
\]
的元素数之和。
类似地，
令
$k_{\mathrm m'}+1$
等于
$\Sigma_{\mathrm m'}$
上的标记点数与集合
\[
\bigcup_{*\ne\mathrm m'}
\bigl(
\Sigma_{\mathrm m'}\cap\Sigma_*
\bigr)
\]
的元素数之和。
最后，令
$k_{\mathrm f,i,j}+1$
等于
$\Sigma_{\mathrm f,i,j}$
上的标记点数与集合
\[
\bigcup_{*\ne(\mathrm f,i,j)}
\bigl(
\Sigma_{\mathrm f,i,j}\cap\Sigma_*
\bigr)
\]
的元素数之和。
于是可以自然地把
$\mathbf S$
视为下列乘积中的元素：
\begin{equation}
\prod_i
\mathcal M^{\mathrm{main}}_{k_{\mathrm m,i}+1}
\times
\prod_{i,j}
\mathcal N_{k_{\mathrm f,i,j}+1}
\times
\mathcal M^{\mathrm{main}}_{k_{\mathrm m'}+1}.
\tag{7.1.57}\label{eq:7-1-57}
\end{equation}
将该元素记为
$\operatorname{dec}(\mathbf S)$。

用恒等于 $0$ 的时间分配，
把
$\mathcal M^{\mathrm{main}}_{k_{\mathrm m,i}+1}$
嵌入
$\mathcal N_{k_{\mathrm m,i}+1}$：
\[
(\Sigma,\vec z)\longmapsto
((\Sigma,\vec z),0).
\]
同样，用恒等于 $1$ 的时间分配，
把
$\mathcal M^{\mathrm{main}}_{k_{\mathrm m'}+1}$
嵌入
$\mathcal N_{k_{\mathrm m'}+1}$。
由此，式~\eqref{eq:7-1-57}
嵌入下列乘积：
\begin{equation}
\prod_i
\mathcal N_{k_{\mathrm m,i}+1}
\times
\prod_{i,j}
\mathcal N_{k_{\mathrm f,i,j}+1}
\times
\mathcal N_{k_{\mathrm m'}+1}.
\tag{7.1.58}\label{eq:7-1-58}
\end{equation}
容易看出，
存在从式~\eqref{eq:7-1-58}
中 $\operatorname{dec}(\mathbf S)$
的邻域
$U(\operatorname{dec}(\mathbf S))$
到
$\mathcal N_{k+1}$
中 $\mathbf S$ 的邻域
$U(\mathbf S)$
的映射，
记为
$\operatorname{dec}_{\mathbf S}^{-1}$。

\begin{statement}{子引理}{7.1.59}\label{sublem:7-1-59}
$\operatorname{dec}_{\mathbf S}^{-1}$
是到一个开子集上的同胚。
\end{statement}

\noindent\textbf{证明。}
关键在于不能把下列分解
\begin{equation}
\Sigma=
\bigcup_i\bigcup_{j=1}^{a_i}
\Sigma_{\mathrm f,i,j}
\cup
\bigcup_i\Sigma_{\mathrm m,i}
\cup\Sigma_{\mathrm m'}
\tag{7.1.60}\label{eq:7-1-60}
\end{equation}
中不同分支交点处的奇点平滑化。
因为若 $p$ 是这些不同分支的交点，
则包含 $p$ 的两个分支有不同的时间分配。
子引理得证。\hfill$\square$

由引理~\ref{lem:7-1-56}
的假设，
分解~\eqref{eq:7-1-60}
不是平凡分解。
因此按归纳假设，
开邻域
$U(\operatorname{dec}(\mathbf S))$
是光滑流形。
引理~\ref{lem:7-1-56} 得证。
\end{proof}

现在证明如下命题。

\begin{proposition}{7.1.61}\label{prop:7-1-61}
$\mathcal N_{k+1}$
具有与按组合型所作分解相容的
带边界或角的光滑流形结构。
\end{proposition}

\begin{proof}
由引理~\ref{lem:7-1-56}，
只须考虑
\[
\mathbf S=(\Sigma,(z_0,\ldots,z_k),\rho)
\]
的
$\overline{\mathbf S}$
仅有一个内部顶点、
且 $\rho$ 为常数的情形。
这样的 $\mathbf S$
的模空间可同
$\mathcal M^{\mathrm{main}}_{k+1}
\times[0,1]$
中的
$\mathfrak I^{-1}(\mathbf S)$
相识别。
第一个因子参数化带标记圆盘的复结构，
第二个因子参数化时间分配 $\rho$。
也就是说，有映射
\begin{equation}
\mathfrak I\colon
\mathcal M^{\mathrm{main}}_{k+1}
\times[0,1]
\longrightarrow
\mathcal N_{k+1},
\tag{7.1.62}\label{eq:7-1-62}
\end{equation}
它给出到所有“带状图仅有一个内部顶点、
且 $\rho$ 为常数”的
$\mathbf S$ 所成集合上的同胚。

若进一步假设 $\Sigma$
只有一个不可约分支，
则
$\mathfrak I^{-1}(\mathbf S)$
属于
$\mathcal M^{\mathrm{main,reg}}_{k+1}
\times[0,1]$。
此外，
由 $\mathfrak I$，
$\mathcal N_{k+1}$
中 $\mathbf S$ 的邻域
同胚于
$\mathcal M^{\mathrm{main,reg}}_{k+1}
\times[0,1]$
中的一个邻域；
因为 $\mathbf S$ 附近的元素
均只有一个分支。
故命题在这样的点附近成立。

现令
$\ell+1$
为 $\Sigma$ 的不可约分支数，
即奇点数为 $\ell$\translatornote{原书首句写
“$\ell+1$ 是奇点数”，
但随即列出 $\ell+1$ 个不可约分支和
$\ell$ 个奇点 $p_i$；
译文按后续分解与角坐标修正，
见疑误 A-0428。}。
此时
$\mathfrak I^{-1}(\mathbf S)$
是
$\mathcal M^{\mathrm{main}}_{k+1}
\times[0,1]$
的角点。
具体地说，
其邻域可与
\[
[0,\epsilon)^\ell
\times
(-\epsilon,\epsilon)^{k-2-\ell}
\times I
\]
相识别；
这里 $I\subset[0,1]$
是 $\rho$ 的开邻域。
下面更精确地说明此同胚。
将 $\Sigma$ 分解为
\[
\Sigma=\bigcup_{i=0}^{\ell}\Sigma_i,
\]
其中 $\Sigma_0$
含第 $0$ 个标记点。
对 $i=1,\ldots,\ell$，
记 $p_i$
为 $\Sigma_i$
同满足
$\Sigma_{j(i)}>\Sigma_i$
的 $\Sigma_{j(i)}$
相交的奇点；
这样的 $p_i$ 唯一。
每个 $\Sigma_i$
上至少有
$k_i+1\ge3$
个奇点或标记点，
它的复结构模空间由
$k_i-2$
个实数参数化，
这些参数刻画这些点的位置。
全部分支合起来给出
$k-2-\ell$
个参数，
构成
$\mathcal M^{\mathrm{main}}_{k+1}$
中
$(-\epsilon,\epsilon)^{k-2-\ell}$
的坐标；
记为
$t_i\in(-\epsilon,\epsilon)$，
$i=1,\ldots,k-2-\ell$。
此外，对
$i=1,\ldots,\ell$，
还有平滑奇点 $p_i$
的参数
$s_i\in[0,\epsilon)$，
它们给出
$[0,\epsilon)^\ell$
因子。

现在讨论时间分配。
设
$\mathbf S'
=(\Sigma',(z_0',\ldots,z_k'),\rho')$
接近 $\mathbf S$。
若奇点 $p_i$
在 $\Sigma'$ 中仍为奇点，
则令
\[
s_i(\mathbf S')
=-\bigl(
\rho'(j(i))-\rho'(i)
\bigr)
\in(-\epsilon,0].
\]
若 $p_i$ 不再是 $\Sigma'$
的奇点，
则上式右侧的两个分支参数没有定义；
此时如上取光滑化参数
$s_i\in(0,\epsilon)$。
同样可以定义
$t_i(\mathbf S')$，
它们参数化
$(-\epsilon,\epsilon)^{k-2-\ell}$。
最后，取
$u(\mathbf S')=\rho'(0)$，
即 $\Sigma_0$ 的时间分配。
于是
\[
(t_1,\ldots,t_{k-2-\ell},
s_1,\ldots,s_\ell,u)
\]
给出从
$\mathbf S$ 邻域到
$(-\epsilon,\epsilon)^{k-2}\times I$
的映射；
其中 $I$ 仍是
$[0,1]$ 中 $\rho$ 的邻域。
容易看出这是同胚。
命题~\ref{prop:7-1-61} 得证。
\end{proof}

边界分解~\eqref{eq:7-1-52-1}、
\eqref{eq:7-1-52-2}
可仿照第~4.6.1 小节的证明；
那里处理的是 $M,L$
不是单点集的情形。
原书特别说，
\eqref{eq:7-1-52-1}
包含 $i=1$、$1+\ell=k$ 的情形，
即
\[
\mathcal M^{\mathrm{main}}_{k+1}
\times\mathcal N_{1+1}
\cong
\mathcal M^{\mathrm{main}}_{k+1},
\]
它对应 $\rho$ 处处为 $1$。
原书还说，
\eqref{eq:7-1-52-2}
包含 $m=k$ 的情形，
即
\[
\mathcal N_{1+1}
\times(\mathcal M^{\mathrm{main}}_{k+1})',
\]
它对应 $\rho$ 处处为 $0$\translatornote{这两个端点例
与定理 7.1.51 给出的指标范围无法逐字相容：
前者在 $i=1,\ell=k-1$ 时应只有
$\ell+1=k$ 个标记点，
后者不满足
$\ell_0=1$ 与严格间隔条件；
详见疑误 A-0429。}。
其余情形是第~4.6.1 小节论证的特例，
故从略。

要完成定理~\ref{thm:7-1-51}
的证明，
还需证明
$\mathcal N_{k+1}$
微分同胚于
$D^{k-1}$；
事实上，本书其他部分并不使用这一事实。
为此构造映射
\[
H\colon
\mathcal N_{k+1}\times[0,1]
\longrightarrow
\mathcal N_{k+1},
\]
\[
H\bigl(
((\Sigma,(z_0,\ldots,z_k)),\rho),
\sigma
\bigr)
=
((\Sigma,(z_0,\ldots,z_k)),\rho_\sigma).
\]
\[
\rho_\sigma(i)
=(1-\sigma)\rho(i)
+\frac{\sigma}{2}.
\]
这里
$\mathbf S=(\Sigma,(z_0,\ldots,z_k),\rho)$，
而原书在给出 $\rho_\sigma$
时写 $\sigma\in(0,1)$。
令 $\mathfrak I$
为式~\eqref{eq:7-1-62} 的映射。

\begin{lemma}{7.1.63}\label{lem:7-1-63}
$H$ 在
$\partial\mathcal N_{k+1}
\times[0,1)$ 上的限制
给出到
\[
\mathcal N_{k+1}\setminus
\mathfrak I\bigl(
\mathcal M^{\mathrm{main}}_{k+1}
\times\{\tfrac12\}
\bigr)
\]
的微分同胚。
\end{lemma}

引理的证明由定义立即可得，
故略去。
回到定理~\ref{thm:7-1-51}
的证明。
$\mathfrak I(
\mathcal M^{\mathrm{main}}_{k+1}
\times\{\tfrac12\})$
微分同胚于
$D^{k-2}$，
并嵌入
$\mathcal N_{k+1}$
的内部；
其正则邻域微分同胚于
$D^{k-1}$。
因此
$\mathcal N_{k+1}$
可缩，
且当 $k\ge3$ 时
$\partial\mathcal N_{k+1}$
单连通。
结论遂由 $h$-配边定理推出。
更具体地说，
$k\ge7$ 时直接由该定理得出；
$k=4,3,2$ 时，
可由维数 $k-1$ 的庞加莱猜想推出；
$k=5,6$ 时，
可以利用引理~7.2.165
的证明将给出的映射
$\mathcal N_{k+1}\to
\mathcal M^{\mathrm{main}}_{k+1}$
的描述，
证明
$\mathcal N_{k+1}\cong D^{k-1}$。
由于本书其他部分并不使用这一事实，
细节从略。
定理~\ref{thm:7-1-51} 得证。\hfill$\square$
